%! Radicals and Rational Exponents — Unit 1, Lesson 1.2 — Exit Ticket
\documentclass[10pt]{article}
\usepackage{atda-article}
\usepackage{atda-boxes}
\newcommand{\TallMath}[1]{$\displaystyle #1\rule[-1.4em]{0pt}{3.2em}$}
\setlength{\workrowsep}{12pt}

\begin{document}

\pageheader{Unit 1, Lesson 1.2}{Exit Ticket}
% NAMESTRIP: no \namedateperiod here. The student writes their name once, on the
% cover this component is stapled behind. See references/conventions.md.

\vspace{0.15in}

\begin{notesbox}{Radicals and Rational Exponents --- On Your Own}
\small
Notes closed. Three items. Leave every answer in exact simplified form.

\begin{enumerate}[leftmargin=1.6em, itemsep=8pt, topsep=6pt]

\item Simplify: $\sqrt{98x^7}$
\begin{work}
  \sqrt{98x^7} &= \sqrt{49x^6 \cdot 2x} \\
               &= 7x^3\sqrt{2x}
\end{work}

\item Combine and simplify: $\sqrt{27}+\sqrt{12}$
\begin{work}
  \sqrt{27}+\sqrt{12} &= 3\sqrt{3}+2\sqrt{3} \\
                      &= 5\sqrt{3}
\end{work}

\item Write $8^{2/3}$ in radical form and evaluate it. Then say in one sentence what the $3$ does
and what the $2$ does.
\begin{work}
  8^{2/3} &= \left(\sqrt[3]{8}\right)^2 \\
          &= 2^2 = 4
\end{work}
\writeline

\end{enumerate}
\end{notesbox}

\end{document}
